Se sabe, por otra parte (\textbf{0}, 23.2.7), que se pueden elegir los
$V_\lambda$ de modo que exista una sub-$A$-álgebra finita $C$ de $K$
tal que $V_\lambda$ sea la clausura integral de
$C_{\mathfrak p_\lambda}$, donde $\mathfrak p_\lambda$ es un ideal
primo de altura 1 en $C$. Si se pone
$C'_{\mathfrak p_\lambda}=C_{\mathfrak p_\lambda}\otimes_A A'$, el
morfismo $\operatorname{Spec}(C'_{\mathfrak p_\lambda})\to
\operatorname{Spec}(C_{\mathfrak p_\lambda})$ es normal (6.8.2) y se
tiene $V'_\lambda=V_\lambda\otimes_{C_{\mathfrak p_\lambda}}
C'_{\mathfrak p_\lambda}$; se puede, pues, reemplazar $B$ por
$V_\lambda$ y $A$ por $C_{\mathfrak p_\lambda}$, y suponer así que
$A$ es local, íntegro y \emph{de dimensión} 1, y que $B$ es su
clausura integral y un anillo de valuación discreta. Existe una
sub-$A$-álgebra \emph{finita} $A_1$ de $B$ tal que el morfismo
$\operatorname{Spec}(B)\to\operatorname{Spec}(A_1)$ es radicial
(\textbf{0}, 23.2.5), lo cual implica en particular que $A_1$ es también un
anillo local (evidentemente de dimensión 1); además, se puede suponer
que $B$ y $A_1$ tienen el mismo cuerpo residual \emph{(loc. cit.)}.
Se puede, pues, por el mismo método, reemplazar $A$ por $A_1$.
Aplicando, si es necesario, el procedimiento inicial de I), se puede
suponer finalmente que $A'$ es también un anillo \emph{local} y que el
homomorfismo $A\to A'$ es \emph{local}.
